%% ma: L12-B08-0005
%% lop: 12
%% chuong: 2
%% bai: 8
%% dang: 12.8.5
%% loai: TLN
%% muc: VD
%% dap_an: "28"
%% nguon: "(NBV)-ĐỀ SỐ 2-MỖI NGÀY 1 ĐỀ THI 2025.pdf (D02, MỖI NGÀY 1 ĐỀ - Đề số 2), Phần III Câu 4 (THPT Văn Giang - Hưng Yên 2025)"
%% kien_thuc: "Bài 7 (tọa độ điểm trong không gian), Bài 8 (ba điểm thẳng hàng, vectơ cùng phương)"
%% ghi_chu: "Đề gốc có hình minh họa (chim, cá, trục tọa độ), không cần vì đã thêm chú thích dấu tọa độ (đề gốc chỉ cho dấu trên hình). Đáp án gốc 28 đúng (kiểm bằng sympy)"
%% trang_thai: cho-duyet
%% ly_do: "Dạng toán mới đề xuất 12.8.5 (tọa độ giao điểm của đường thẳng đi qua hai điểm với mặt phẳng Oxy bằng điều kiện thẳng hàng), chờ giáo viên duyệt."
\begin{ex}
Với hệ trục tọa độ $Oxyz$ sao cho $O$ nằm trên mặt nước, mặt phẳng $(Oxy)$ là mặt nước, trục $Oz$ hướng lên trên (đơn vị đo: mét), một con chim bói cá đang ở vị trí cách mặt nước $2$ m, cách mặt phẳng $(Oxz)$, $(Oyz)$ lần lượt là $3$ m và $1$ m phóng thẳng xuống vị trí con cá cách mặt nước $50$ cm, cách mặt phẳng $(Oxz)$, $(Oyz)$ lần lượt là $1$ m và $1{,}5$ m (cả hai vị trí đều có hoành độ và tung độ dương). Tọa độ điểm $B$ lúc chim bói cá vừa tiếp xúc với mặt nước là $(a;b;c)$. Tính $T=5a+15b+25c$.
\shortans{28}
\loigiai{Chọn đơn vị $1$ m. Vị trí chim là $C(1;3;2)$, vị trí cá là $A(1{,}5;1;-0{,}5)$. Điểm $B$ nằm trên mặt nước nên $c=0$ và $B$ thuộc đường thẳng $CA$, tức $\overrightarrow{CB}=k\overrightarrow{CA}$:
$(a-1;\,b-3;\,-2)=k(0{,}5;\,-2;\,-2{,}5)\Rightarrow k=\dfrac45$, $a=\dfrac75$, $b=\dfrac75$.
Vậy $T=5\cdot\dfrac75+15\cdot\dfrac75+25\cdot0=28$.}
\end{ex}
